\chapter{The Growing Tree}
\label{ch:growing_tree}

The spine of the guide is an eight-iteration search on our four-legal-move endgame position. Because there are only four legal moves at the root, every single node in the tree can be drawn explicitly at fixed coordinates without schematic shortcuts.

\section{Selection Rule Recall}

Before tracing the search, we recall the core score calculation derived in §1.4 (Equation~\eqref{eq:E10}).

\begin{established}[Selection Rule Recall]
Selection score $S(a) = Q(a) + U(a)$ (Equation~\eqref{eq:E10}) governs move choice (see \textbf{FIG-1.10} in §1.5 for the complete term-by-term derivation and engine parameters). Unvisited moves inherit $Q_{\text{FPU}}$ floor (Equation~\eqref{eq:E11}) penalized by FPU reduction. The move with the tallest bar ($\max S$) is selected.
\end{established}

\clearpage

\begin{realdata}[Where the Solid Bar Segment Comes From]
The origin of solid bar segments before and after measurement:
\begin{center}
\begin{minipage}{\linewidth}
\captionsetup{type=figure}
\centering
\subcaptionbox{1. $V$ arrives}{% GENERATED by tools/make_figures.py — do not edit.
\begin{tikzpicture}
  \node[vgroot] (root) at (0,0) {Root\\ \scriptsize $V(\mathrm{root}) = +0.97602$};
  \node[vgunvisited, inner sep=1.8pt] (kd6) at (-4.5,-1.8) {Kd6\\ \scriptsize unvisited};
  \node[vgunvisited, inner sep=1.8pt] (kf6) at (-1.5,-1.8) {Kf6\\ \scriptsize unvisited};
  \node[vgunvisited, inner sep=1.8pt] (kf5) at (1.5,-1.8) {Kf5\\ \scriptsize unvisited};
  \node[vgunvisited, inner sep=1.8pt] (kd5) at (4.5,-1.8) {Kd5\\ \scriptsize unvisited};
  \draw[vgedge] (root) -- (kd6);
  \draw[vgedge] (root) -- (kf6);
  \draw[vgedge] (root) -- (kf5);
  \draw[vgedge] (root) -- (kd5);
  \begin{scope}[yshift=-6.8cm]
    \node[anchor=west, font=\scriptsize\itshape, text=WorkGrey] at (-6.0, 3.4) {Step 1: $V$ arrives at root. Move bars are empty:};
    \vgsbaraxis{-6.0}{6.0}{1.8}
    \draw[NavyBlue!40, dash pattern=on 2pt off 1.5pt] (-4.75,0) rectangle (-4.25, 1.763);
    \draw[NavyBlue!40, dash pattern=on 2pt off 1.5pt] (-1.75,0) rectangle (-1.25, 1.748);
    \draw[NavyBlue!40, dash pattern=on 2pt off 1.5pt] (1.25,0) rectangle (1.75, 1.070);
    \draw[NavyBlue!40, dash pattern=on 2pt off 1.5pt] (4.25,0) rectangle (4.75, 1.068);
  \end{scope}
\end{tikzpicture}
}\qquad
\subcaptionbox{2. $V$ becomes $Q_{\mathrm{FPU}}$ floor}{% GENERATED by tools/make_figures.py — do not edit.
\begin{tikzpicture}
  \node[vgroot] (root) at (0,0) {Root\\ \scriptsize $V(\mathrm{root}) = +0.97602$};
  \node[vgunvisited, inner sep=1.8pt] (kd6) at (-4.5,-1.8) {Kd6\\ \scriptsize $Q_{\mathrm{FPU}}{=}0.97602$};
  \node[vgunvisited, inner sep=1.8pt] (kf6) at (-1.5,-1.8) {Kf6\\ \scriptsize $Q_{\mathrm{FPU}}{=}0.97602$};
  \node[vgunvisited, inner sep=1.8pt] (kf5) at (1.5,-1.8) {Kf5\\ \scriptsize $Q_{\mathrm{FPU}}{=}0.97602$};
  \node[vgunvisited, inner sep=1.8pt] (kd5) at (4.5,-1.8) {Kd5\\ \scriptsize $Q_{\mathrm{FPU}}{=}0.97602$};
  \draw[vgedge] (root) -- (kd6);
  \draw[vgedge] (root) -- (kf6);
  \draw[vgedge] (root) -- (kf5);
  \draw[vgedge] (root) -- (kd5);
  \begin{scope}[yshift=-6.8cm]
    \node[anchor=west, font=\scriptsize\itshape, text=WorkGrey] at (-6.0, 3.4) {Step 2: $V$ sets $Q_{\mathrm{FPU}}$ floor for ALL moves:};
    \vgsbaraxis{-6.0}{6.0}{1.8}
    \vgsbarS{-4.5}{0.97602}{0.78756}{Kd6}{fpu}
    \vgsbarS{-1.5}{0.97602}{0.77186}{Kf6}{fpu}
    \vgsbarS{1.5}{0.97602}{0.09389}{Kf5}{fpu}
    \vgsbarS{4.5}{0.97602}{0.09179}{Kd5}{fpu}
  \end{scope}
\end{tikzpicture}
}\\[2ex]
\subcaptionbox{3. FPU penalty appears}{% GENERATED by tools/make_figures.py — do not edit.
\begin{tikzpicture}
  \node[vgroot] (root) at (0,0) {Root\\ \scriptsize $n{=}1$\quad $Q{=}0.97184$};
  \node[vgnode, fill=BuildBlue!40, inner sep=1.8pt] (kd6) at (-4.5,-1.8) {Kd6\\ \scriptsize $n{=}1$\ \ $Q{=}0.96766$};
  \node[vgunvisited, inner sep=1.8pt] (kf6) at (-1.5,-1.8) {Kf6\\ \scriptsize $Q_{\mathrm{FPU}}{=}0.75015$};
  \node[vgunvisited, inner sep=1.8pt] (kf5) at (1.5,-1.8) {Kf5\\ \scriptsize $Q_{\mathrm{FPU}}{=}0.75015$};
  \node[vgunvisited, inner sep=1.8pt] (kd5) at (4.5,-1.8) {Kd5\\ \scriptsize $Q_{\mathrm{FPU}}{=}0.75015$};
  \vgedgewidth{1}{1}\draw[vgedge, line width=\vgEW pt] (root) -- (kd6);
  \vgedgewidth{0}{1}\draw[vgedge, line width=\vgEW pt] (root) -- (kf6);
  \vgedgewidth{0}{1}\draw[vgedge, line width=\vgEW pt] (root) -- (kf5);
  \vgedgewidth{0}{1}\draw[vgedge, line width=\vgEW pt] (root) -- (kd5);
  \begin{scope}[yshift=-6.8cm]
    \node[anchor=west, font=\scriptsize\itshape, text=WorkGrey] at (-6.0, 3.4) {Step 3: After Kd6 visit, penalty $0.33\sqrt{0.4513} = 0.22169$ drops unvisited bars:};
    \vgsbaraxis{-6.0}{6.0}{1.8}
    \vgsbarS{-4.5}{0.96766}{0.39378}{Kd6}{meas}
    \vgsbarS{-1.5}{0.75015}{0.77186}{Kf6}{fpu}
    \vgsbarS{1.5}{0.75015}{0.09389}{Kf5}{fpu}
    \vgsbarS{4.5}{0.75015}{0.09179}{Kd5}{fpu}

    \draw[PitfallRed, fill=PitfallRed!25, line width=0.4pt, dash pattern=on 1pt off 1.5pt] (-1.75, 1.20) rectangle (-1.25, 1.56);
    \node[font=\tiny\bfseries, text=PitfallRed] at (-1.5, 1.38) {$-0.222$};
  \end{scope}

  \node[draw=NavyBlue, fill=white, rounded corners=3pt, inner sep=4pt, align=center, below=52mm of root, font=\scriptsize] (table) {
    \textbf{Running $Q_{\mathrm{FPU}}$ Tracking Across 8 Iterations}\\[0.5ex]
    \begin{tabular}{ccc}
    \textbf{Iteration} & \textbf{Visited Prior Mass} & \textbf{$Q_{\mathrm{FPU}}$} \\
    \hline
    1 & 0.0000 & 0.97602 \\
    2 & 0.4513 & 0.75015 \\
    3 & 0.8936 & 0.66460 \\
    4--8 & 0.8936 & 0.658 -- 0.668 \\
    \end{tabular}\\[0.5ex]
    \textit{The penalty stops growing after iteration 3 because only Kd6 and Kf6 are ever visited,}\\
    \textit{so no new prior mass is added. $Q_{\mathrm{FPU}}$ only drifts with root $Q$ (see \texttt{ch07:fpudrop}).}
  };
\end{tikzpicture}
}
\captionof{figure}{\textbf{FIG-2.B: Where the solid bar segment comes from.} (a) $V(\text{root}) = +0.97602$ arrives (Equation~\eqref{eq:E2}). (b) Before any move is visited, $V$ sets the equal baseline floor ($Q_{\text{FPU}} = +0.97602$, Equation~\eqref{eq:E11}) for all 4 moves; $P$ sets the pale $U$ caps. (c) After Kd6 is measured, a penalty slice ($0.33\sqrt{0.4513} = 0.22169$) drops unvisited FPU scores to $0.75015$.}
\label{fig:2_b}
\end{minipage}
\end{center}
\end{realdata}

\clearpage

\section{The Eight Search Iterations}

We trace the search iteration by iteration. A FIG-2.x frame captures two moments: the \textbf{tree} displays the state \textit{after} iteration $x$ completes (new leaf evaluated, visit counts incremented, backpropagated values), while the \textbf{S-bar strip} below displays the scores compared \textit{before} iteration $x$'s pick --- the numbers that caused the selection. The tallest bar ($S = Q + U$) carries gold highlight and matches the candidate move marked with the gold selection ring in the tree above.

\begin{realdata}[Iteration 0]
\begin{center}
% GENERATED by tools/make_figures.py — do not edit.
\begin{tikzpicture}
  \node[vgroot] (root) at (0,0) {Root\\ \scriptsize $n{=}0$\quad $Q{=}0.97602$};
  \node[vgunvisited, inner sep=1.8pt] (kd6) at (-4.5,-1.8) {Kd6\\ \scriptsize $n{=}0$\ \ $Q_{\mathrm{FPU}}{=}0.97602$};
  \node[vgunvisited, inner sep=1.8pt] (kf6) at (-1.5,-1.8) {Kf6\\ \scriptsize $n{=}0$\ \ $Q_{\mathrm{FPU}}{=}0.97602$};
  \node[vgunvisited, inner sep=1.8pt] (kf5) at (1.5,-1.8) {Kf5\\ \scriptsize $n{=}0$\ \ $Q_{\mathrm{FPU}}{=}0.97602$};
  \node[vgunvisited, inner sep=1.8pt] (kd5) at (4.5,-1.8) {Kd5\\ \scriptsize $n{=}0$\ \ $Q_{\mathrm{FPU}}{=}0.97602$};
  \vgedgewidth{0}{1}\draw[vgedge, line width=\vgEW pt] (root) -- (kd6);
  \vgedgewidth{0}{1}\draw[vgedge, line width=\vgEW pt] (root) -- (kf6);
  \vgedgewidth{0}{1}\draw[vgedge, line width=\vgEW pt] (root) -- (kf5);
  \vgedgewidth{0}{1}\draw[vgedge, line width=\vgEW pt] (root) -- (kd5);
  \begin{scope}[yshift=-9.5cm]
    \node[anchor=west, font=\scriptsize\itshape, text=WorkGrey] at (-6.0, 3.6) {Initial priors and FPU scores before Iteration 1 ($S = Q_{\mathrm{FPU}} + U$):};
    \vgsbaraxis{-6.0}{6.0}{1.8}
    \vgsbarS{-4.5}{0.97602}{0.78756}{Kd6}{fpu}
    \vgsbarS{-1.5}{0.97602}{0.77186}{Kf6}{fpu}
    \vgsbarS{1.5}{0.97602}{0.09389}{Kf5}{fpu}
    \vgsbarS{4.5}{0.97602}{0.09179}{Kd5}{fpu}
  \end{scope}
\end{tikzpicture}

\captionof{figure}{\textbf{FIG-2.0: Iteration 0 (Starting State).} All four children unvisited (dashed grey). All $Q$ values are equal to $Q_{\text{FPU}} = +0.97602$ (Equation~\eqref{eq:E11}). Ranking is purely decided by policy priors.}
\label{fig:2_0}
\end{center}
\end{realdata}

\vspace{1ex}

\begin{realdata}[Iteration 1]
\begin{center}
% GENERATED by tools/make_figures.py — do not edit.

\begin{tikzpicture}
  \node[vgroot] (root) at (0,0) {Root\\ \scriptsize $n{=}1$\quad $Q{=}+0.97184$};
  \vgheat{1}{1}\node[vgnew, fill=vgfill, inner sep=1.8pt] (kd6) at (-4.5,-1.8) {Kd6\\ \scriptsize $n{=}1$\ \ $Q{=}+0.96766$};
  \node[vgunvisited, inner sep=1.8pt] (kf6) at (-1.5,-1.8) {Kf6\\ \scriptsize $n{=}0$\ \ $Q_{\mathrm{FPU}}{=}0.75015$};
  \node[vgunvisited, inner sep=1.8pt] (kf5) at (1.5,-1.8) {Kf5\\ \scriptsize $n{=}0$\ \ $Q_{\mathrm{FPU}}{=}0.75015$};
  \node[vgunvisited, inner sep=1.8pt] (kd5) at (4.5,-1.8) {Kd5\\ \scriptsize $n{=}0$\ \ $Q_{\mathrm{FPU}}{=}0.75015$};
  \vgedgewidth{1}{1}\draw[vgedge, line width=\vgEW pt] (root) -- (kd6);
  \vgedgewidth{0}{1}\draw[vgedge, line width=\vgEW pt] (root) -- (kf6);
  \vgedgewidth{0}{1}\draw[vgedge, line width=\vgEW pt] (root) -- (kf5);
  \vgedgewidth{0}{1}\draw[vgedge, line width=\vgEW pt] (root) -- (kd5);
  \vgselectmark{kd6}
  \begin{scope}[yshift=-9.5cm]
    \node[anchor=west, font=\scriptsize\itshape, text=WorkGrey] at (-6.0, 3.6) {Scores compared BEFORE Iteration 1 pick ($S = Q + U$):};
    \vgsbaraxis{-6.0}{6.0}{1.8}
    \vgsbarS{-4.5}{0.97602}{0.78756}{Kd6}{sel}
    \vgsbarS{-1.5}{0.97602}{0.77186}{Kf6}{fpu}
    \vgsbarS{1.5}{0.97602}{0.09389}{Kf5}{fpu}
    \vgsbarS{4.5}{0.97602}{0.09179}{Kd5}{fpu}
  \end{scope}
\end{tikzpicture}

\captionof{figure}{\textbf{FIG-2.1: Iteration 1 (Selection $\rightarrow$ Expansion $\rightarrow$ Evaluation $\rightarrow$ Backpropagation).} Kd6 selected ($S = 1.76358$, Equation~\eqref{eq:E10}). Expanded leaf returns $+0.96766$. Kd6 becomes solid green burst; backpropagation updates root to $Q = +0.97184$ (Equation~\eqref{eq:E4}). Search stopped at depth 1.}
\label{fig:2_1}
\end{center}
\end{realdata}
\clearpage

\begin{realdata}[Iteration 2]
\begin{center}
% GENERATED by tools/make_figures.py — do not edit.

\begin{tikzpicture}
  \node[vgroot] (root) at (0,0) {Root\\ \scriptsize $n{=}2$\quad $Q{=}+0.97655$};
  \vgheat{1}{1}\node[vgnode, fill=vgfill, inner sep=1.8pt] (kd6) at (-4.5,-1.8) {Kd6\\ \scriptsize $n{=}1$\ \ $Q{=}+0.96766$};
  \vgheat{1}{1}\node[vgnew, fill=vgfill, inner sep=1.8pt] (kf6) at (-1.5,-1.8) {Kf6\\ \scriptsize $n{=}1$\ \ $Q{=}+0.98598$};
  \node[vgunvisited, inner sep=1.8pt] (kf5) at (1.5,-1.8) {Kf5\\ \scriptsize $n{=}0$\ \ $Q_{\mathrm{FPU}}{=}0.66460$};
  \node[vgunvisited, inner sep=1.8pt] (kd5) at (4.5,-1.8) {Kd5\\ \scriptsize $n{=}0$\ \ $Q_{\mathrm{FPU}}{=}0.66460$};
  \vgedgewidth{1}{1}\draw[vgedge, line width=\vgEW pt] (root) -- (kd6);
  \vgedgewidth{1}{1}\draw[vgedge, line width=\vgEW pt] (root) -- (kf6);
  \vgedgewidth{0}{1}\draw[vgedge, line width=\vgEW pt] (root) -- (kf5);
  \vgedgewidth{0}{1}\draw[vgedge, line width=\vgEW pt] (root) -- (kd5);
  \vgselectmark{kf6}
  \begin{scope}[yshift=-9.5cm]
    \node[anchor=west, font=\scriptsize\itshape, text=WorkGrey] at (-6.0, 3.6) {Scores compared BEFORE Iteration 2 pick ($S = Q + U$):};
    \vgsbaraxis{-6.0}{6.0}{1.8}
    \vgsbarS{-4.5}{0.96766}{0.39380}{Kd6}{meas}
    \vgsbarS{-1.5}{0.75015}{0.77190}{Kf6}{sel}
    \vgsbarS{1.5}{0.75015}{0.09389}{Kf5}{fpu}
    \vgsbarS{4.5}{0.75015}{0.09180}{Kd5}{fpu}
  \end{scope}
\end{tikzpicture}

\captionof{figure}{\textbf{FIG-2.2: Iteration 2 (Selection $\rightarrow$ Expansion $\rightarrow$ Evaluation $\rightarrow$ Backpropagation).} Kd6 has the best measured score ($+0.96766$ vs pre-pick FPU $0.75015$, grey bars) yet is \textbf{not} selected. Its $U$ halved ($1+0 \rightarrow 1+1$ denominator in Equation~\eqref{eq:E10}), allowing Kf6 ($S = 1.52205$) to take the lead. Post-pick FPU updates to $0.66460$ (Equation~\eqref{eq:E11}).}
\label{fig:2_2}
\end{center}
\end{realdata}

\vspace{1ex}

\begin{realdata}[Iteration 3]
\begin{center}
% GENERATED by tools/make_figures.py — do not edit.

\begin{tikzpicture}
  \node[vgroot] (root) at (0,0) {Root\\ \scriptsize $n{=}3$\quad $Q{=}+0.97024$};
  \vgheat{1}{2}\node[vgnode, fill=vgfill, inner sep=1.8pt] (kd6) at (-4.5,-1.8) {Kd6\\ \scriptsize $n{=}1$\ \ $Q{=}+0.96766$};
  \vgheat{2}{2}\node[vgnode, fill=vgfill, inner sep=1.8pt] (kf6) at (-1.5,-1.8) {Kf6\\ \scriptsize $n{=}2$\ \ $Q{=}+0.96863$};
  \node[vgunvisited, inner sep=1.8pt] (kf5) at (1.5,-1.8) {Kf5\\ \scriptsize $n{=}0$\ \ $Q_{\mathrm{FPU}}{=}0.65829$};
  \node[vgunvisited, inner sep=1.8pt] (kd5) at (4.5,-1.8) {Kd5\\ \scriptsize $n{=}0$\ \ $Q_{\mathrm{FPU}}{=}0.65829$};
  \vgedgewidth{1}{2}\draw[vgedge, line width=\vgEW pt] (root) -- (kd6);
  \vgedgewidth{2}{2}\draw[vgedge, line width=\vgEW pt] (root) -- (kf6);
  \vgedgewidth{0}{2}\draw[vgedge, line width=\vgEW pt] (root) -- (kf5);
  \vgedgewidth{0}{2}\draw[vgedge, line width=\vgEW pt] (root) -- (kd5);
  \node[vgnew] (kf8) at (-1.0,-3.6) {Kf8\\ \scriptsize leaf $V{=}+0.95129$};
  \draw[vgedge] (kf6) -- (kf8);
  \draw[vgpath] (root) -- (kf6);
  \draw[vgpath] (kf6) -- (kf8);
  \vgnewmark{kf8}
  \vgselectmark{kf6}
  \vglane{kf8}{root}{6.5}{-4.5}{$+0.95129$}
  \vgselectmark{kf6}
  \begin{scope}[yshift=-9.5cm]
    \node[anchor=west, font=\scriptsize\itshape, text=WorkGrey] at (-6.0, 3.6) {Scores compared BEFORE Iteration 3 pick ($S = Q + U$):};
    \vgsbaraxis{-6.0}{6.0}{1.8}
    \vgsbarS{-4.5}{0.96766}{0.55696}{Kd6}{meas}
    \vgsbarS{-1.5}{0.98598}{0.54585}{Kf6}{sel}
    \vgsbarS{1.5}{0.66460}{0.13279}{Kf5}{fpu}
    \vgsbarS{4.5}{0.66460}{0.12983}{Kd5}{fpu}
  \end{scope}
\end{tikzpicture}

\captionof{figure}{\textbf{FIG-2.3: Iteration 3 (Selection recursed twice before expanding).} Kf6 selected ($S = 1.53183$). Selection does not stop at root; it recurses down Kf6 to expand leaf Kf8 (depth 2). Leaf value $+0.95129$ backs up.}
\label{fig:2_3}
\end{center}
\end{realdata}
\clearpage

\begin{realdata}[Iteration 4]
\begin{center}
% GENERATED by tools/make_figures.py — do not edit.

\begin{tikzpicture}
  \node[vgroot] (root) at (0,0) {Root\\ \scriptsize $n{=}4$\quad $Q{=}+0.97571$};
  \vgheat{2}{2}\node[vgnode, fill=vgfill, inner sep=1.8pt] (kd6) at (-4.5,-1.8) {Kd6\\ \scriptsize $n{=}2$\ \ $Q{=}+0.98262$};
  \vgheat{2}{2}\node[vgnode, fill=vgfill, inner sep=1.8pt] (kf6) at (-1.5,-1.8) {Kf6\\ \scriptsize $n{=}2$\ \ $Q{=}+0.96863$};
  \node[vgunvisited, inner sep=1.8pt] (kf5) at (1.5,-1.8) {Kf5\\ \scriptsize $n{=}0$\ \ $Q_{\mathrm{FPU}}{=}0.66376$};
  \node[vgunvisited, inner sep=1.8pt] (kd5) at (4.5,-1.8) {Kd5\\ \scriptsize $n{=}0$\ \ $Q_{\mathrm{FPU}}{=}0.66376$};
  \vgedgewidth{2}{2}\draw[vgedge, line width=\vgEW pt] (root) -- (kd6);
  \vgedgewidth{2}{2}\draw[vgedge, line width=\vgEW pt] (root) -- (kf6);
  \vgedgewidth{0}{2}\draw[vgedge, line width=\vgEW pt] (root) -- (kf5);
  \vgedgewidth{0}{2}\draw[vgedge, line width=\vgEW pt] (root) -- (kd5);
  \node[vgnode] (kf8) at (-1.0,-3.6) {Kf8\\ \scriptsize leaf $V{=}+0.95129$};
  \draw[vgedge] (kf6) -- (kf8);
  \node[vgnew] (kd8) at (-5.8,-3.6) {Kd8\\ \scriptsize leaf $V{=}+0.99759$};
  \draw[vgedge] (kd6) -- (kd8);
  \draw[vgpath] (root) -- (kd6);
  \draw[vgpath] (kd6) -- (kd8);
  \vgnewmark{kd8}
  \vgselectmark{kd6}
  \vgselectmark{kd6}
  \begin{scope}[yshift=-9.5cm]
    \node[anchor=west, font=\scriptsize\itshape, text=WorkGrey] at (-6.0, 3.6) {Scores compared BEFORE Iteration 4 pick ($S = Q + U$):};
    \vgsbaraxis{-6.0}{6.0}{1.8}
    \vgsbarS{-4.5}{0.96766}{0.68217}{Kd6}{sel}
    \vgsbarS{-1.5}{0.96863}{0.44571}{Kf6}{meas}
    \vgsbarS{1.5}{0.65829}{0.16264}{Kf5}{fpu}
    \vgsbarS{4.5}{0.65829}{0.15902}{Kd5}{fpu}
  \end{scope}
\end{tikzpicture}

\captionof{figure}{\textbf{FIG-2.4: Iteration 4 (Selection recursed twice before expanding).} Kd6 selected ($S = 1.64983$). Selection recurses down Kd6 branch to expand new leaf Kd8 (depth 2). Leaf returns $+0.99759$, updating Kd6's average $Q$ to $+0.98262$.}
\label{fig:2_4}
\end{center}
\end{realdata}

\vspace{1ex}

\begin{realdata}[Iteration 5]
\begin{center}
% GENERATED by tools/make_figures.py — do not edit.

\begin{tikzpicture}
  \node[vgroot] (root) at (0,0) {Root\\ \scriptsize $n{=}5$\quad $Q{=}+0.97974$};
  \vgheat{3}{3}\node[vgnode, fill=vgfill, inner sep=1.8pt] (kd6) at (-4.5,-1.8) {Kd6\\ \scriptsize $n{=}3$\ \ $Q{=}+0.98839$};
  \vgheat{2}{3}\node[vgnode, fill=vgfill, inner sep=1.8pt] (kf6) at (-1.5,-1.8) {Kf6\\ \scriptsize $n{=}2$\ \ $Q{=}+0.96863$};
  \node[vgunvisited, inner sep=1.8pt] (kf5) at (1.5,-1.8) {Kf5\\ \scriptsize $n{=}0$\ \ $Q_{\mathrm{FPU}}{=}0.66779$};
  \node[vgunvisited, inner sep=1.8pt] (kd5) at (4.5,-1.8) {Kd5\\ \scriptsize $n{=}0$\ \ $Q_{\mathrm{FPU}}{=}0.66779$};
  \vgedgewidth{3}{3}\draw[vgedge, line width=\vgEW pt] (root) -- (kd6);
  \vgedgewidth{2}{3}\draw[vgedge, line width=\vgEW pt] (root) -- (kf6);
  \vgedgewidth{0}{3}\draw[vgedge, line width=\vgEW pt] (root) -- (kf5);
  \vgedgewidth{0}{3}\draw[vgedge, line width=\vgEW pt] (root) -- (kd5);
  \node[vgnode] (kf8) at (-1.0,-3.6) {Kf8\\ \scriptsize leaf $V{=}+0.95129$};
  \draw[vgedge] (kf6) -- (kf8);
  \node[vgnode] (kd8) at (-5.8,-3.6) {Kd8\\ \scriptsize leaf $V{=}+0.99759$};
  \draw[vgedge] (kd6) -- (kd8);
  \node[vgnew] (kf7) at (-3.4,-3.6) {Kf7\\ \scriptsize leaf $V{=}+0.99992$};
  \draw[vgedge] (kd6) -- (kf7);
  \draw[vgpath] (root) -- (kd6);
  \draw[vgpath] (kd6) -- (kf7);
  \vgnewmark{kf7}
  \vgselectmark{kd6}
  \vglane{kf7}{root}{-6.5}{-4.5}{$+0.99992$}
  \vgselectmark{kd6}
  \begin{scope}[yshift=-9.5cm]
    \node[anchor=west, font=\scriptsize\itshape, text=WorkGrey] at (-6.0, 3.6) {Scores compared BEFORE Iteration 5 pick ($S = Q + U$):};
    \vgsbaraxis{-6.0}{6.0}{1.8}
    \vgsbarS{-4.5}{0.98262}{0.52516}{Kd6}{sel}
    \vgsbarS{-1.5}{0.96863}{0.51469}{Kf6}{meas}
    \vgsbarS{1.5}{0.66376}{0.18782}{Kf5}{fpu}
    \vgsbarS{4.5}{0.66376}{0.18363}{Kd5}{fpu}
  \end{scope}
\end{tikzpicture}

\captionof{figure}{\textbf{FIG-2.5: Iteration 5 (Selection recursed twice before expanding).} Kd6 selected ($S = 1.50779$). Selection recurses down Kd6 to expand Kf7 (depth 2). Leaf returns $+0.99992$. Kd6 subtree widens.}
\label{fig:2_5}
\end{center}
\end{realdata}
\clearpage

\begin{realdata}[Iteration 6]
\begin{center}
% GENERATED by tools/make_figures.py — do not edit.

\begin{tikzpicture}
  \node[vgroot] (root) at (0,0) {Root\\ \scriptsize $n{=}6$\quad $Q{=}+0.97958$};
  \vgheat{3}{3}\node[vgnode, fill=vgfill, inner sep=1.8pt] (kd6) at (-4.5,-1.8) {Kd6\\ \scriptsize $n{=}3$\ \ $Q{=}+0.98839$};
  \vgheat{3}{3}\node[vgnode, fill=vgfill, inner sep=1.8pt] (kf6) at (-1.5,-1.8) {Kf6\\ \scriptsize $n{=}3$\ \ $Q{=}+0.97196$};
  \node[vgunvisited, inner sep=1.8pt] (kf5) at (1.5,-1.8) {Kf5\\ \scriptsize $n{=}0$\ \ $Q_{\mathrm{FPU}}{=}0.66763$};
  \node[vgunvisited, inner sep=1.8pt] (kd5) at (4.5,-1.8) {Kd5\\ \scriptsize $n{=}0$\ \ $Q_{\mathrm{FPU}}{=}0.66763$};
  \vgedgewidth{3}{3}\draw[vgedge, line width=\vgEW pt] (root) -- (kd6);
  \vgedgewidth{3}{3}\draw[vgedge, line width=\vgEW pt] (root) -- (kf6);
  \vgedgewidth{0}{3}\draw[vgedge, line width=\vgEW pt] (root) -- (kf5);
  \vgedgewidth{0}{3}\draw[vgedge, line width=\vgEW pt] (root) -- (kd5);
  \node[vgnode] (kf8) at (-1.0,-3.6) {Kf8\\ \scriptsize leaf $V{=}+0.95129$};
  \draw[vgedge] (kf6) -- (kf8);
  \node[vgnode] (kd8) at (-5.8,-3.6) {Kd8\\ \scriptsize leaf $V{=}+0.99759$};
  \draw[vgedge] (kd6) -- (kd8);
  \node[vgnode] (kf7) at (-3.4,-3.6) {Kf7\\ \scriptsize leaf $V{=}+0.99992$};
  \draw[vgedge] (kd6) -- (kf7);
  \node[vgnew] (e6_kf8) at (-1.0,-5.2) {e6\\ \scriptsize leaf $V{=}+0.98598$};
  \draw[vgedge] (kf8) -- (e6_kf8);
  \draw[vgpath] (root) -- (kf6);
  \draw[vgpath] (kf6) -- (kf8);
  \draw[vgpath] (kf8) -- (e6_kf8);
  \vgnewmark{e6_kf8}
  \vgselectmark{kf8}
  \vglane{e6_kf8}{root}{6.5}{-6.0}{$+0.98598$}
  \vgselectmark{kf6}
  \begin{scope}[yshift=-9.5cm]
    \node[anchor=west, font=\scriptsize\itshape, text=WorkGrey] at (-6.0, 3.6) {Scores compared BEFORE Iteration 6 pick ($S = Q + U$):};
    \vgsbaraxis{-6.0}{6.0}{1.8}
    \vgsbarS{-4.5}{0.98839}{0.44039}{Kd6}{meas}
    \vgsbarS{-1.5}{0.96863}{0.57547}{Kf6}{sel}
    \vgsbarS{1.5}{0.66779}{0.21000}{Kf5}{fpu}
    \vgsbarS{4.5}{0.66779}{0.20531}{Kd5}{fpu}
  \end{scope}
\end{tikzpicture}

\captionof{figure}{\textbf{FIG-2.6: Iteration 6 (Selection recursed three times before expanding).} Kf6 selected ($S = 1.54411$). Selection reaches depth 3 at e6 (under Kf8). Leaf returns $+0.97860$. Search deepens autonomously.}
\label{fig:2_6}
\end{center}
\end{realdata}

\vspace{1ex}

\begin{realdata}[Iteration 7]
\begin{center}
% GENERATED by tools/make_figures.py — do not edit.

\begin{tikzpicture}
  \node[vgroot] (root) at (0,0) {Root\\ \scriptsize $n{=}7$\quad $Q{=}+0.97846$};
  \vgheat{4}{4}\node[vgnode, fill=vgfill, inner sep=1.8pt] (kd6) at (-4.5,-1.8) {Kd6\\ \scriptsize $n{=}4$\ \ $Q{=}+0.98394$};
  \vgheat{3}{4}\node[vgnode, fill=vgfill, inner sep=1.8pt] (kf6) at (-1.5,-1.8) {Kf6\\ \scriptsize $n{=}3$\ \ $Q{=}+0.97196$};
  \node[vgunvisited, inner sep=1.8pt] (kf5) at (1.5,-1.8) {Kf5\\ \scriptsize $n{=}0$\ \ $Q_{\mathrm{FPU}}{=}0.66651$};
  \node[vgunvisited, inner sep=1.8pt] (kd5) at (4.5,-1.8) {Kd5\\ \scriptsize $n{=}0$\ \ $Q_{\mathrm{FPU}}{=}0.66651$};
  \vgedgewidth{4}{4}\draw[vgedge, line width=\vgEW pt] (root) -- (kd6);
  \vgedgewidth{3}{4}\draw[vgedge, line width=\vgEW pt] (root) -- (kf6);
  \vgedgewidth{0}{4}\draw[vgedge, line width=\vgEW pt] (root) -- (kf5);
  \vgedgewidth{0}{4}\draw[vgedge, line width=\vgEW pt] (root) -- (kd5);
  \node[vgnode] (kf8) at (-1.0,-3.6) {Kf8\\ \scriptsize leaf $V{=}+0.95129$};
  \draw[vgedge] (kf6) -- (kf8);
  \node[vgnode] (kd8) at (-5.8,-3.6) {Kd8\\ \scriptsize leaf $V{=}+0.99759$};
  \draw[vgedge] (kd6) -- (kd8);
  \node[vgnode] (kf7) at (-3.4,-3.6) {Kf7\\ \scriptsize leaf $V{=}+0.99992$};
  \draw[vgedge] (kd6) -- (kf7);
  \node[vgnode] (e6_kf8) at (-1.0,-5.2) {e6\\ \scriptsize leaf $V{=}+0.98598$};
  \draw[vgedge] (kf8) -- (e6_kf8);
  \node[vgnew] (e6_kd8) at (-5.8,-5.2) {e6\\ \scriptsize leaf $V{=}+0.98598$};
  \draw[vgedge] (kd8) -- (e6_kd8);
  \draw[vgpath] (root) -- (kd6);
  \draw[vgpath] (kd6) -- (kd8);
  \draw[vgpath] (kd8) -- (e6_kd8);
  \vgnewmark{e6_kd8}
  \vgselectmark{kd8}
  \vglane{e6_kd8}{root}{-6.5}{-6.0}{$+0.98598$}
  \vgselectmark{kd6}
  \begin{scope}[yshift=-9.5cm]
    \node[anchor=west, font=\scriptsize\itshape, text=WorkGrey] at (-6.0, 3.6) {Scores compared BEFORE Iteration 7 pick ($S = Q + U$):};
    \vgsbaraxis{-6.0}{6.0}{1.8}
    \vgsbarS{-4.5}{0.98839}{0.48245}{Kd6}{sel}
    \vgsbarS{-1.5}{0.97196}{0.47283}{Kf6}{meas}
    \vgsbarS{1.5}{0.66763}{0.23005}{Kf5}{fpu}
    \vgsbarS{4.5}{0.66763}{0.22492}{Kd5}{fpu}
  \end{scope}
\end{tikzpicture}

\captionof{figure}{\textbf{FIG-2.7: Iteration 7 (Selection recursed three times before expanding).} Kd6 selected ($S = 1.47084$). Selection reaches depth 3 at e6 (under Kd8). Leaf returns $+0.97060$.}
\label{fig:2_7}
\end{center}
\end{realdata}
\clearpage

\begin{realdata}[Iteration 8]
\begin{center}
\input{figures/fig_2_8.tex}
\captionof{figure}{\textbf{FIG-2.8: Iteration 8 (8 full MCTS cycles complete).} State after 7 backed-up values. Hand arithmetic matches \texttt{lc0.exe go nodes 8} visit counts and $Q$ values to 5 decimal places.}
\label{fig:2_8}
\end{center}
\end{realdata}

\vspace{1ex}

\begin{realdata}[Search Time-Lapse]
\begin{center}
% GENERATED by tools/make_figures.py — do not edit.
\begin{tikzpicture}[scale=0.85, every node/.style={transform shape}]
  % Row 1: Iterations 1-4
  \begin{scope}[xshift=0cm, yshift=0cm]
    \node[vgroot, minimum width=8mm, minimum height=3mm] (r1) at (0,0) {};
    \node[vgnode, fill=BuildBlue!40, minimum width=6mm, minimum height=3mm] (kd6_1) at (-1.2,-1.0) {};
    \node[vgunvisited, minimum width=6mm, minimum height=3mm] (kf6_1) at (-0.4,-1.0) {};
    \node[vgunvisited, minimum width=6mm, minimum height=3mm] (kf5_1) at (0.4,-1.0) {};
    \node[vgunvisited, minimum width=6mm, minimum height=3mm] (kd5_1) at (1.2,-1.0) {};
    \draw[vgedge, line width=1.5pt] (r1) -- (kd6_1);
    \draw[vgedge, line width=0.4pt] (r1) -- (kf6_1);
    \draw[vgedge, line width=0.4pt] (r1) -- (kf5_1);
    \draw[vgedge, line width=0.4pt] (r1) -- (kd5_1);
    \node[above=1mm of r1, font=\scriptsize\bfseries] {It. 1};
  \end{scope}

  \begin{scope}[xshift=4.5cm, yshift=0cm]
    \node[vgroot, minimum width=8mm, minimum height=3mm] (r2) at (0,0) {};
    \node[vgnode, fill=BuildBlue!40, minimum width=6mm, minimum height=3mm] (kd6_2) at (-1.2,-1.0) {};
    \node[vgnode, fill=BuildBlue!40, minimum width=6mm, minimum height=3mm] (kf6_2) at (-0.4,-1.0) {};
    \node[vgunvisited, minimum width=6mm, minimum height=3mm] (kf5_2) at (0.4,-1.0) {};
    \node[vgunvisited, minimum width=6mm, minimum height=3mm] (kd5_2) at (1.2,-1.0) {};
    \draw[vgedge, line width=1.5pt] (r2) -- (kd6_2);
    \draw[vgedge, line width=1.5pt] (r2) -- (kf6_2);
    \draw[vgedge, line width=0.4pt] (r2) -- (kf5_2);
    \draw[vgedge, line width=0.4pt] (r2) -- (kd5_2);
    \node[above=1mm of r2, font=\scriptsize\bfseries] {It. 2};
  \end{scope}

  \begin{scope}[xshift=9.0cm, yshift=0cm]
    \node[vgroot, minimum width=8mm, minimum height=3mm] (r3) at (0,0) {};
    \node[vgnode, fill=BuildBlue!40, minimum width=6mm, minimum height=3mm] (kd6_3) at (-1.2,-1.0) {};
    \node[vgnode, fill=BuildBlue!50, minimum width=6mm, minimum height=3mm] (kf6_3) at (-0.4,-1.0) {};
    \node[vgunvisited, minimum width=6mm, minimum height=3mm] (kf5_3) at (0.4,-1.0) {};
    \node[vgunvisited, minimum width=6mm, minimum height=3mm] (kd5_3) at (1.2,-1.0) {};
    \node[vgnode, fill=BuildBlue!40, minimum width=6mm, minimum height=3mm] (kf8_3) at (-0.4,-2.0) {};
    \draw[vgedge, line width=1.5pt] (r3) -- (kd6_3);
    \draw[vgedge, line width=2.0pt] (r3) -- (kf6_3);
    \draw[vgedge, line width=1.5pt] (kf6_3) -- (kf8_3);
    \draw[vgedge, line width=0.4pt] (r3) -- (kf5_3);
    \draw[vgedge, line width=0.4pt] (r3) -- (kd5_3);
    \node[above=1mm of r3, font=\scriptsize\bfseries] {It. 3};
  \end{scope}

  \begin{scope}[xshift=13.5cm, yshift=0cm]
    \node[vgroot, minimum width=8mm, minimum height=3mm] (r4) at (0,0) {};
    \node[vgnode, fill=BuildBlue!50, minimum width=6mm, minimum height=3mm] (kd6_4) at (-1.2,-1.0) {};
    \node[vgnode, fill=BuildBlue!50, minimum width=6mm, minimum height=3mm] (kf6_4) at (-0.4,-1.0) {};
    \node[vgunvisited, minimum width=6mm, minimum height=3mm] (kf5_4) at (0.4,-1.0) {};
    \node[vgunvisited, minimum width=6mm, minimum height=3mm] (kd5_4) at (1.2,-1.0) {};
    \node[vgnode, fill=BuildBlue!40, minimum width=6mm, minimum height=3mm] (kd8_4) at (-1.6,-2.0) {};
    \node[vgnode, fill=BuildBlue!40, minimum width=6mm, minimum height=3mm] (kf8_4) at (-0.4,-2.0) {};
    \draw[vgedge, line width=2.0pt] (r4) -- (kd6_4);
    \draw[vgedge, line width=1.5pt] (kd6_4) -- (kd8_4);
    \draw[vgedge, line width=2.0pt] (r4) -- (kf6_4);
    \draw[vgedge, line width=1.5pt] (kf6_4) -- (kf8_4);
    \draw[vgedge, line width=0.4pt] (r4) -- (kf5_4);
    \draw[vgedge, line width=0.4pt] (r4) -- (kd5_4);
    \node[above=1mm of r4, font=\scriptsize\bfseries] {It. 4};
  \end{scope}

  % Row 2: Iterations 5-8
  \begin{scope}[xshift=0cm, yshift=-3.6cm]
    \node[vgroot, minimum width=8mm, minimum height=3mm] (r5) at (0,0) {};
    \node[vgnode, fill=BuildBlue!60, minimum width=6mm, minimum height=3mm] (kd6_5) at (-1.2,-1.0) {};
    \node[vgnode, fill=BuildBlue!50, minimum width=6mm, minimum height=3mm] (kf6_5) at (-0.4,-1.0) {};
    \node[vgunvisited, minimum width=6mm, minimum height=3mm] (kf5_5) at (0.4,-1.0) {};
    \node[vgunvisited, minimum width=6mm, minimum height=3mm] (kd5_5) at (1.2,-1.0) {};
    \node[vgnode, fill=BuildBlue!40, minimum width=6mm, minimum height=3mm] (kd8_5) at (-1.6,-2.0) {};
    \node[vgnode, fill=BuildBlue!40, minimum width=6mm, minimum height=3mm] (kf7_5) at (-0.9,-2.0) {};
    \node[vgnode, fill=BuildBlue!40, minimum width=6mm, minimum height=3mm] (kf8_5) at (-0.2,-2.0) {};
    \draw[vgedge, line width=2.2pt] (r5) -- (kd6_5);
    \draw[vgedge, line width=1.5pt] (kd6_5) -- (kd8_5);
    \draw[vgedge, line width=1.5pt] (kd6_5) -- (kf7_5);
    \draw[vgedge, line width=2.0pt] (r5) -- (kf6_5);
    \draw[vgedge, line width=1.5pt] (kf6_5) -- (kf8_5);
    \draw[vgedge, line width=0.4pt] (r5) -- (kf5_5);
    \draw[vgedge, line width=0.4pt] (r5) -- (kd5_5);
    \node[above=1mm of r5, font=\scriptsize\bfseries] {It. 5};
  \end{scope}

  \begin{scope}[xshift=4.5cm, yshift=-3.6cm]
    \node[vgroot, minimum width=8mm, minimum height=3mm] (r6) at (0,0) {};
    \node[vgnode, fill=BuildBlue!60, minimum width=6mm, minimum height=3mm] (kd6_6) at (-1.2,-1.0) {};
    \node[vgnode, fill=BuildBlue!60, minimum width=6mm, minimum height=3mm] (kf6_6) at (-0.4,-1.0) {};
    \node[vgunvisited, minimum width=6mm, minimum height=3mm] (kf5_6) at (0.4,-1.0) {};
    \node[vgunvisited, minimum width=6mm, minimum height=3mm] (kd5_6) at (1.2,-1.0) {};
    \node[vgnode, fill=BuildBlue!40, minimum width=6mm, minimum height=3mm] (kd8_6) at (-1.6,-2.0) {};
    \node[vgnode, fill=BuildBlue!40, minimum width=6mm, minimum height=3mm] (kf7_6) at (-0.9,-2.0) {};
    \node[vgnode, fill=BuildBlue!45, minimum width=6mm, minimum height=3mm] (kf8_6) at (-0.2,-2.0) {};
    \node[vgnode, fill=BuildBlue!40, minimum width=6mm, minimum height=3mm] (e6_kf8_6) at (-0.2,-3.0) {};
    \draw[vgedge, line width=2.2pt] (r6) -- (kd6_6);
    \draw[vgedge, line width=1.5pt] (kd6_6) -- (kd8_6);
    \draw[vgedge, line width=1.5pt] (kd6_6) -- (kf7_6);
    \draw[vgedge, line width=2.2pt] (r6) -- (kf6_6);
    \draw[vgedge, line width=1.8pt] (kf6_6) -- (kf8_6);
    \draw[vgedge, line width=1.5pt] (kf8_6) -- (e6_kf8_6);
    \draw[vgedge, line width=0.4pt] (r6) -- (kf5_6);
    \draw[vgedge, line width=0.4pt] (r6) -- (kd5_6);
    \node[above=1mm of r6, font=\scriptsize\bfseries] {It. 6};
  \end{scope}

  \begin{scope}[xshift=9.0cm, yshift=-3.6cm]
    \node[vgroot, minimum width=8mm, minimum height=3mm] (r7) at (0,0) {};
    \node[vgnode, fill=BuildBlue!70, minimum width=6mm, minimum height=3mm] (kd6_7) at (-1.2,-1.0) {};
    \node[vgnode, fill=BuildBlue!60, minimum width=6mm, minimum height=3mm] (kf6_7) at (-0.4,-1.0) {};
    \node[vgunvisited, minimum width=6mm, minimum height=3mm] (kf5_7) at (0.4,-1.0) {};
    \node[vgunvisited, minimum width=6mm, minimum height=3mm] (kd5_7) at (1.2,-1.0) {};
    \node[vgnode, fill=BuildBlue!45, minimum width=6mm, minimum height=3mm] (kd8_7) at (-1.6,-2.0) {};
    \node[vgnode, fill=BuildBlue!40, minimum width=6mm, minimum height=3mm] (kf7_7) at (-0.9,-2.0) {};
    \node[vgnode, fill=BuildBlue!45, minimum width=6mm, minimum height=3mm] (kf8_7) at (-0.2,-2.0) {};
    \node[vgnode, fill=BuildBlue!40, minimum width=6mm, minimum height=3mm] (e6_kd8_7) at (-1.6,-3.0) {};
    \node[vgnode, fill=BuildBlue!40, minimum width=6mm, minimum height=3mm] (e6_kf8_7) at (-0.2,-3.0) {};
    \draw[vgedge, line width=2.5pt] (r7) -- (kd6_7);
    \draw[vgedge, line width=1.8pt] (kd6_7) -- (kd8_7);
    \draw[vgedge, line width=1.5pt] (kd6_7) -- (kf7_7);
    \draw[vgedge, line width=1.5pt] (kd8_7) -- (e6_kd8_7);
    \draw[vgedge, line width=2.2pt] (r7) -- (kf6_7);
    \draw[vgedge, line width=1.8pt] (kf6_7) -- (kf8_7);
    \draw[vgedge, line width=1.5pt] (kf8_7) -- (e6_kf8_7);
    \draw[vgedge, line width=0.4pt] (r7) -- (kf5_7);
    \draw[vgedge, line width=0.4pt] (r7) -- (kd5_7);
    \node[above=1mm of r7, font=\scriptsize\bfseries] {It. 7};
  \end{scope}

  \begin{scope}[xshift=13.5cm, yshift=-3.6cm]
    \node[vgroot, minimum width=8mm, minimum height=3mm] (r8) at (0,0) {};
    \node[vgnode, fill=BuildBlue!70, minimum width=6mm, minimum height=3mm] (kd6_8) at (-1.2,-1.0) {};
    \node[vgnode, fill=BuildBlue!65, minimum width=6mm, minimum height=3mm] (kf6_8) at (-0.4,-1.0) {};
    \node[vgunvisited, minimum width=6mm, minimum height=3mm] (kf5_8) at (0.4,-1.0) {};
    \node[vgunvisited, minimum width=6mm, minimum height=3mm] (kd5_8) at (1.2,-1.0) {};
    \node[vgnode, fill=BuildBlue!45, minimum width=6mm, minimum height=3mm] (kd8_8) at (-1.6,-2.0) {};
    \node[vgnode, fill=BuildBlue!40, minimum width=6mm, minimum height=3mm] (kf7_8) at (-0.9,-2.0) {};
    \node[vgnode, fill=BuildBlue!45, minimum width=6mm, minimum height=3mm] (kf8_8) at (-0.2,-2.0) {};
    \node[vgnode, fill=BuildBlue!40, minimum width=6mm, minimum height=3mm] (e6_kd8_8) at (-1.6,-3.0) {};
    \node[vgnode, fill=BuildBlue!40, minimum width=6mm, minimum height=3mm] (e6_kf8_8) at (-0.2,-3.0) {};
    \draw[vgedge, line width=2.5pt] (r8) -- (kd6_8);
    \draw[vgedge, line width=1.8pt] (kd6_8) -- (kd8_8);
    \draw[vgedge, line width=1.5pt] (kd6_8) -- (kf7_8);
    \draw[vgedge, line width=1.5pt] (kd8_8) -- (e6_kd8_8);
    \draw[vgedge, line width=2.4pt] (r8) -- (kf6_8);
    \draw[vgedge, line width=1.8pt] (kf6_8) -- (kf8_8);
    \draw[vgedge, line width=1.5pt] (kf8_8) -- (e6_kf8_8);
    \draw[vgedge, line width=0.4pt] (r8) -- (kf5_8);
    \draw[vgedge, line width=0.4pt] (r8) -- (kd5_8);
    \node[above=1mm of r8, font=\scriptsize\bfseries] {It. 8};
  \end{scope}
\end{tikzpicture}

\captionof{figure}{\textbf{FIG-2.9: The whole search time-lapse.} Sequence of tree states across iterations 1 to 8 showing structural growth.}
\label{fig:2_9}
\end{center}
\end{realdata}
\clearpage

\section{The Sign Flip Convention}

Evaluations are recorded from the point of view of the player to move at that node. As values walk back up the tree, the sign flips at every ply (Equation~\eqref{eq:E8}):

\begin{realdata}[Sign Flips]
\begin{center}
\begin{minipage}{\linewidth}
\captionsetup{type=figure}
\centering
\subcaptionbox{Leaf Evaluation (White POV)}{% GENERATED by tools/make_figures.py — do not edit.
\begin{tikzpicture}
  \vgboardnode{b1}{(0,0)}{5k2/8/5K2/4P3/8/8/8/8 w - - 2 2}{11pt}
  \node[draw=IdeaGreen, fill=IdeaGreen!10, rounded corners=3pt, inner sep=6pt, right=8mm of b1, text width=52mm, font=\small] (info) {
    \textbf{Leaf Node (after 1.Kf6 Kf8)}\\[0.5ex]
    Side to move: \textbf{\textcolor{BuildBlue}{White}}\\[0.5ex]
    Evaluated $V = \mathbf{+0.95129}$\\[0.5ex]
    \textit{"White has a winning advantage (+0.951)."}
  };
\end{tikzpicture}
}\\[1ex]
\subcaptionbox{Black-to-move Edge}{% GENERATED by tools/make_figures.py — do not edit.
\begin{tikzpicture}
  \node[vgnode, draw=NavyBlue] (parent) at (0,0) {Kf6 Node\\ \scriptsize Black to move};
  \node[vgnew, right=22mm of parent] (leaf) {Kf8 Leaf\\ \scriptsize White to move};
  \draw[vgbackprop] (leaf) to node[midway, above, font=\scriptsize\bfseries, text=PitfallRed, fill=white] {$-0.95129$} (parent);
  \node[draw=PitfallRed, fill=PitfallRed!8, rounded corners=3pt, inner sep=6pt, below=4mm of parent, text width=60mm, font=\small] (note) {
    \textbf{Incoming Edge to Kf8 (Black to move)}\\[0.5ex]
    Black evaluates White's $+0.95129$ win as $\mathbf{-0.95129}$ for Black.
  };
\end{tikzpicture}
}\\[1ex]
\subcaptionbox{Root Edge (White POV)}{% GENERATED by tools/make_figures.py — do not edit.
\begin{tikzpicture}
  \node[vgroot] (root) at (0,0) {Root\\ \scriptsize White to move};
  \node[vgnode, right=22mm of root] (kf6) {Kf6 Node\\ \scriptsize Black to move};
  \node[vgnew, right=22mm of kf6] (leaf) {Kf8 Leaf\\ \scriptsize White to move};
  \draw[vgbackprop, bend right=35] (leaf) to (kf6);
  \draw[vgbackprop, bend right=35] (kf6) to node[midway, above, font=\scriptsize\bfseries, text=IdeaGreen, fill=white] {$+0.95129$} (root);
  \node[draw=IdeaGreen, fill=IdeaGreen!8, rounded corners=3pt, inner sep=6pt, below=4mm of root, text width=70mm, font=\small] (note) {
    \textbf{Root Edge (White to move)}\\[0.5ex]
    Sign flips back: White receives $\mathbf{+0.95129}$.\\[0.3ex]
    Kf6 running average: $(0.98598 + 0.95129) / 2 = \mathbf{0.968635}$.
  };
\end{tikzpicture}
}
\captionof{figure}{\textbf{FIG-2.10: Iteration 3 Sign Flips.} New leaf after 1.Kf6 Kf8 gives $+0.95129$ (White). Black edge receives $-0.95129$ via Negamax sign flip (Equation~\eqref{eq:E8}). Root edge receives $+0.95129$, updating Kf6 average to $0.968635$.}
\label{fig:2_10}
\end{minipage}
\end{center}
\end{realdata}

\section{Emergent Depth}

\begin{established}[Autonomous Search Expansion]
\begin{center}
% GENERATED by tools/make_figures.py — do not edit.
\begin{tikzpicture}
  \draw[draw=NavyBlue, fill=PaleGrey, rounded corners=3pt, inner sep=6pt] (0,0) rectangle (10.5,2.8);
  \node[anchor=north west, font=\small\bfseries, text=NavyBlue] at (0.3,2.6) {Depth Emergence across Iterations};
  \node[anchor=north west, font=\small] at (0.3,2.0) {
    \begin{tabular}{@{}ll@{}}
      Iterations 1--2: & Stop at root's children (\textbf{depth 1})\\[0.3ex]
      Iterations 3--5: & Reach \textbf{depth 2} (Kf8, Kd8, Kf7 appear)\\[0.3ex]
      Iterations 6--7: & Reach \textbf{depth 3} (e6 pawn moves appear)\\[0.3ex]
    \end{tabular}
  };
  \node[anchor=south west, font=\small\itshape, text=IdeaGreen] at (0.3,0.2) {
    Leela has no depth parameter at all. Depth is what happens when $U$ collapses.
  };
\end{tikzpicture}

\captionof{figure}{\textbf{FIG-2.11: Depth Emergence.} Iterations 1--2 stop at depth 1; 3--5 reach depth 2; 6--7 reach depth 3. Depth is an emergent outcome of $U$ decay (Equation~\eqref{eq:E10}), not a pre-configured parameter.}
\label{fig:2_11}
\end{center}
\end{established}

